\subsection{\identifier{UC-10} --- Record Attendance and Actual Work Periods}
\label{uc:10}

\paragraph{Purpose}

Allow an authorised user to confirm daily attendance
and record actual work periods, tasks, locations
and optional breaks.

Planned information may support entry,
but becomes actual work only after
explicit confirmation.


\paragraph{Actors}

\textbf{Primary actor:} Farm Owner.

\textbf{Additional actor:} Field Manager,
with the relevant permissions granted by the owner.

The authorised owner or manager records worker attendance.
Workers do not need application accounts for this recording.

\paragraph{Preconditions}

\begin{itemize}
    \item The user has permission
    for the requested operation.
    \item Worker records exist centrally or were created locally through \identifier{UC-07}.
    A new worker may be created before returning to attendance/work entry.
    \item The relevant campaign and work areas
    exist within the authorised scope.
    \item Work dates comply with
    the campaign lifecycle rules.
    \item Offline entry requires the necessary
    authorised reference data to be available.
\end{itemize}

\paragraph{Main Flow}

\begin{enumerate}
    \item The user selects the campaign
    and working date.
    \item The system displays planned workers,
    available proposed periods and any existing
    attendance and work records.
    \item Only selected/expected workers without a daily decision are shown
    as awaiting confirmation. Unrelated unselected workers generate no alert.
    \item The user marks workers as present or absent.
    Daily attendance is shared within the business
    for that worker and working date.
    \item For present workers, the user identifies
    the actual campaign and work-area assignments.
    These may remain incomplete as to hours or tasks.
    \item The user reviews proposed periods
    or enters actual periods, tasks and areas.
    Common values may be applied to a group
    with individual adjustments.
    \item The user reviews or records actual breaks,
    including their start, end and whether
    they count towards working time.
    Planned breaks are not accepted automatically.
    \item The system validates the entries
    and calculates available working-time totals.
    \item The user reviews and explicitly
    confirms the submitted information.
    \item The system saves the records
    and displays their local storage
    and synchronisation status.
\end{enumerate}

\paragraph{Alternative Flows and Exceptions}

\begin{description}

    \item[A1 --- Attendance is known before hours.]
    The user saves presence or absence
    and completes work details later.
    Missing hours remain incomplete,
    not a confirmed zero total.

    \item[A2 --- An unplanned worker attends.]
    The user selects an existing authorised
    worker record, including a worker created offline with authorisation through \identifier{UC-07},
    and records attendance.
    A planned assignment is not required.

    \item[A3 --- No plan exists.]
    The user selects existing workers
    and enters the actual information directly.

    \item[A4 --- Actual work differs from the plan.]
    The user adjusts workers, periods,
    tasks or areas before confirmation.
    The original plan remains unchanged.
    A complete period requires its actual task before confirmation.
    Until then, preserve the incomplete assignment rather than inventing hours.

    \item[A5 --- A worker changes task or area.]
    The user records separate periods for the relevant campaign
    and area. Each confirmed complete period identifies its actual task;
    attendance/assignment may be saved earlier without completed periods.

    \item[A6 --- Work periods exclude a break.]
    The interval between those periods
    does not contribute to working time.
    The same excluded interval is not
    deducted again.

    \item[A7 --- A break occurs within a period.]
    The user records its start and end.
    It is deducted only when explicitly
    marked as not counted towards working time.

    \item[A8 --- A worker is absent.]
    The system does not convert planned
    periods into actual work.
    Contradictory actual work entries
    require resolution.

    \item[A9 --- Work crosses midnight.]
    The period retains its start-date association.
    Its end date and time establish
    the full interval for calculations
    and overlap validation.

    \item[A10 --- Work is entered retrospectively.]
    The user selects the actual working date.
    The system preserves the distinction
    between work date and entry date.

    \item[A11 --- The campaign is closed.]

    The user may complete previously missing information
    for work performed within the campaign's valid actual dates.

    This includes entering missing details or a previously
    unrecorded actual period, provided that confirmed values
    and associations remain unchanged.

    The system validates attendance consistency, permissions,
    campaign coverage and time overlaps.

    Changing or removing confirmed information requires
    authorised reopening and follows \identifier{UC-11}.

    Work outside the campaign's valid actual dates requires
    reopening and date review.

    \item[A12 --- Confirmed information needs correction.]
    The user follows \identifier{UC-11}.
    Existing values are not silently replaced
    through another initial submission.

    \item[A13 --- Invalid or overlapping periods.]
    The system identifies the problem
    and preserves unaffected input.
    Validation considers the worker's relevant
    periods across areas and campaigns.

    \item[A14 --- The device is offline.]
    The system saves supported entries locally.
    Their synchronisation status remains visible.

    Recording actual work does not silently change campaign dates.

    If the campaign has no actual start, the first confirmed
    actual assignment provides a proposed start date under \identifier{UC-06}.

    The user may explicitly confirm the proposed start offline. Start and
    dependent work remain pending until the server validates them in order;
    a competing campaign start is preserved for resolution, not silently accepted.
\end{description}

\paragraph{Postconditions}

\textbf{On success:}
The submitted attendance and available actual
work details are saved in the correct context.

\textbf{On attendance-only confirmation:}
Presence is recorded while hours and tasks
remain explicitly incomplete where unknown.

\textbf{On offline save:}
User-confirmed entries are stored locally
and remain pending server acknowledgement.

\textbf{On cancellation or rejected changes:}
Previously saved information remains unchanged.

\textbf{In all cases:}
Actual work entry does not rewrite the plan.

Actual-period validation follows \identifier{FR-033} and Section~\ref{sec:counted-working-time}.

\paragraph{Related Requirements and Rules}
\hyperref[fr:004]{\identifier{FR-004}}, \hyperref[fr:016]{\identifier{FR-016}}, \hyperref[fr:020]{\identifier{FR-020}}, \hyperref[fr:030]{\identifier{FR-030}}, \hyperref[fr:031]{\identifier{FR-031}}, \hyperref[fr:032]{\identifier{FR-032}}, \hyperref[fr:033]{\identifier{FR-033}}, \hyperref[fr:034]{\identifier{FR-034}}, \hyperref[fr:035]{\identifier{FR-035}}.

\hyperref[br:011]{\identifier{BR-011}}, \hyperref[br:018]{\identifier{BR-018}}, \hyperref[br:019]{\identifier{BR-019}}, \hyperref[br:020]{\identifier{BR-020}}, \hyperref[br:021]{\identifier{BR-021}}.

\paragraph{Decision References}
\decisionref{Q-WF-01}, \decisionref{Q-WF-03}, \decisionref{D-04}.
